\cotefile{12}
\scanpage{36}
% L1
induisant la filtration donnée, \deleted{\unsure{On a}}
\deleted{ici}
\[
M\in\mathrm{Ob}\,D_i(M(X))\Longleftrightarrow M(-j)\in\mathrm{Ob}\,D_i(M(X))\quad\text{pour $j$ \unsure{grand|grd}}
\]
On a
\[
\deleted{D_i^{\unsure{-}}(M^+(X))/D_{i-1}(\unsure{\bar M|M|\mathcal M}^+(X))\xrightarrow{\sim}D_i(M(X))/D_{i-1}(M(X))}
\]
\[
D_i(M(X))\otimes D_j(M(X))\subset D_{i+j}(M(X))
\]
et de $=$ en mettant des $+$. \inserted{mais $\mathbb Z(-j)\in\mathrm{Ob}\,D_0(M(X))$ pour tt $j$.}

% L7
Enfin, on a
\[
\underline{\underline{\mathrm{Ext}}}^{\deleted{\unsure{k}}i}(D_j(M(X)),D_k(M(X))\subset D_{\deleted{\unsure{i}}-j+k}(M(X))
\]
(qui se formule aussi simple en termes des $M_{\alpha}(X)$ !).

% L11
10) \uline{Motifs constants tordus}. \uline{Anneaux} $\unsure{\mathbf M|\mathbb M|M}^+(X)$ et $\mathbf M(X)$

Soit $M^+(X)_{\mathrm{spec.}}$ la ss-catégorie des des motifs effectifs \uline{constants} \uline{tordus}.
[\deleted{\unsure{que|Par}} condition locale, qui se vérifie sur $T_\ell(M)$, cf N\textsuperscript{o}~8], et de $=$ on définit $M(X)_{\mathrm{spec.}}$ \deleted{\unsure{Notons|N.A.} que $M$ est} Ainsi $M$ est tordu $\Longleftrightarrow M(j)$ est tordu.

% L18
$M^+(X)_{\mathrm{spec}}$ et $M(X)_{\mathrm{spec}}$ sont des \inserted{ss-}\uline{catégories abéliennes}, \uline{stables par extensions}, \uline{dont tt objet est de}
% archive sheet number: 35
\scanpage{37}
% L1
\uuline{longueur finie}. Elle est stable par $\otimes$, sur cette catégorie les $\mathrm{Ext}^i$ $(i>0)$ sont nuls, et la catégorie est aussi stable par $\mathscr{H}\mathrm{om}$. (cas de $M(X)_{\mathrm{spec}}$). On a une ``catégorie tensorielle \unsure{''} sur $\mathbb Q$''.

% L6
Le \deleted{\unsure{la}} $\lambda$-anneau défini par cette catégorie est noté $\unsure{\mathbf M|\mathbb M|M}^+(X)$ resp. $M(X)$. Noter que $\underline{\xi}=\mathrm{classe}\,\deleted{\mathbb Z}\mathbb Q_X(-1))$ est un élément de $M^+(X)$, et on a
\[
M(X)=M^+(X)_\xi=M^+(X)\left[\frac{1}{\xi}\right].
\]

% L11
En fait, $\xi$ est non div. de $0$ : $\mathbf M^+(X)$ comme $\mathbb Z$-module admet comme base les classes d'\deleted{\unsure{éléments}}\inserted{\uline{objets}} \uline{simples} de $M^+(X)$, et $x\mapsto\xi\cdot x$ dans $\mathbf M^+(X)$ est défini par une application \uline{injective} sur l'ens de ses objets simples. [N.B. $M$ simple \inserted{ds $\mathbf M^+(X)$} $\Longleftrightarrow$ \inserted{$M$ simple dans $M(X)$ $\Longleftrightarrow$} $\mathbf M(-1)$ simple, comme il résulte du fait que $\mathbf M^+(X)$ est épaisse dans $M(X)$ et $\otimes\mathbb Z(-1)$ une auto équivalence ]
% archive sheet number: 36
\scanpage{38}
% L1
Si $\Sigma^+(X)=$ ens des classes d'objets simples de $M^+(X)$, $\Sigma(X)=\underset{\substack{\text{par}\\e\mapsto\xi e}}{\underrightarrow{\lim}}\Sigma^+(X)$, on aura \deleted{\unsure{pour}} \deleted{$M(K)$} un diagramme comm.
\diagram{\begin{tikzcd}[ampersand replacement=\&,column sep=small,row sep=large]
M^+(X)\arrow[d,hook] \& = \& \mathbb Z^{(\Sigma^+(X))}\arrow[d,hook]\\
M(X) \& = \& \mathbb Z^{(\Sigma(X)}
\end{tikzcd}}

% L6
Alors $M_i^{\unsure{+}}(X)$ est défini en termes d'une partie $\Sigma^+_{(i)}(X)$ de $\Sigma_i(X)$, d'où
\[
M^+_i(X)=\mathbb Z^{\Sigma^+_{(i)}(X)}
\]
ce qui donne une filtration croissante sur l'anneau $M^+(X)$, et une filtration analogue sur $M(X)$.

% L12
En fait, posant
\[
\Sigma_i^+(X)=\Sigma^+_{(i)}(X)-\Sigma^+_{(\unsure{i-1|i+1})}(X)
\]
on aura $=$
\[
\Sigma_i^+(X)\cdot\Sigma_j^+(X)\subset\mathbb Z^{\Sigma^+_{\unsure{i+j|i-j}}(X)}
\]
donc $M^+(X)$ devient un \uline{anneau}
% archive sheet number: 37
\scanpage{39}
% L1
\uline{gradué} \inserted{\uline{à degrés} $\geq 0$}, et $\xi$ est de degré $1$. De cette façon, $M(X)$ devient gradué [à degrés des deux signes, $\xi$ étant un élément \uline{inversible} \unsure{en|de} degré $1$]. Noter que
\[
M_0(X)=\mathbb Z^{\Sigma_0(X)}
\]
avec
\[
\Sigma_0(X)=\underset{\text{par }x\mapsto x\cdot\xi}{\underrightarrow{\lim}}\Sigma_i^+(X)
\]
et
\[
\begin{aligned}
M^{\mathrm{pair}}(X)&=\underbrace{M_0(X)}_{\text{de degré }0}[\underbrace{\xi}_{\text{de degré }1}]\\
M^{\mathrm{impair}}(X)&=M_1(X)\otimes_{M_0(X)}M_0(X)[\xi]
\end{aligned}
\]
% L10
donc la structure d'anneau gradué de $M(X)$ est déterminée en termes de la structure d'anneau de $M_0(X)$, \deleted{et} le Module $M_1(X)$ sur $M_0(X)$, et l'application de mult.
$M_1(X)\otimes_{M_{0}(X)}M_1(X)\longrightarrow M_0(X)$ \unsure{définie par} \deleted{\unsure{$\xi$}} $u,v\mapsto\xi^{-1}uv$.

% L15
D'autre part, les objets de $D_i^+(M^+(X))$ sont \deleted{\unsure{précist}} \inserted{ceux} dont la classe dans \deleted{$\unsure{M|K}^+(X)$ sont} $M^+(X)$ sont dans $\mathbb Z^{D_i(\Sigma^+(X))}$, où
\[
\left.\begin{aligned}
D_i(\Sigma^+(X))&=\coprod_{j\geq0}\deleted{\unsure{\mathbb Z^{x}}}\xi^j\Sigma^+_{(i)}(X)\\
&=\coprod_{\substack{j\geq0\\0\leq k\leq i}}\xi^j\Sigma_k^+(X)
\end{aligned}\right\}\unsure{\underset{\mathrm{df}}{=}|\mathrel{=}|\xi^j}\Sigma^+_{[j]}(X)
\]
% archive sheet number: 38
\scanpage{40}
% L1
Posons
\[
\deleted{\Sigma}\deleted{\unsure{\Sigma}}\sigma_i(X)=\Sigma_i^+(X)-\bigcup_{\substack{\deleted{\unsure{0<j<i}}\\\alpha\geq0}}\xi^{\alpha\deleted{\unsure{j}}}\Sigma^+_{\deleted{\unsure{j}}i-2\alpha}(X)
\]
Alors on a une partition (au sens large)
\[
\deleted{\unsure{\Sigma}}\Sigma^{\unsure{+}}_i(X)=\sigma_i(X)\cup\xi\sigma_{i-2}(X)\cup\dotsb\cup
\begin{cases}\xi^{i/2}\sigma_0(X)\\ \xi^{\frac{i-1}{2}}\sigma_1(X)\end{cases}
\]
% margin note written sideways
\margin{\uline{suivant la parité}}
% slanted
\margin{éléments \uline{primitifs} de poids $i$}
et on a donc
\[
D_i(\Sigma^+(X))\mathbin{\deleted{\unsure{=}}\cap}\Sigma_j^+(X)=
\begin{cases}
\Sigma_j^+(X)&\text{si }j\leq i\\
\multicolumn{2}{l}{\deleted{\unsure{e}}\xi^{j-i}\sigma_i(X)\cup\xi^{j-i-1}\sigma_{i-1}(X)\cup}\\
\multicolumn{2}{r}{\dotsb\cup\xi^j\sigma_0(X)}\\
&\text{si }j>i
\end{cases}
\]

% L8
Ainsi
\[
\begin{aligned}
D_i(M^+(X))/D_{i-1}(M^+(X))&\simeq\mathbb Z^{D_i(\Sigma^+(X))-D_{i-1}(\Sigma^+(X))}\\
&\simeq\mathbb Z^{\sigma_i(X)+\xi\sigma_{i-2}(X)+\xi^{2}\sigma_{i-4}(X)+\dotsb,}\\
&\deleted{\unsure{\simeq}\mathbb Z^{\sigma_{\unsure{i}}(X)}}
\end{aligned}
\]

% L12
\uline{N.B.} Du pt de vue de la structure additive, des deux filtrations, et de l'opération $\xi\cdot$, $M^+(X)$ et $M(X)$ sont complètement déterminés par la connaissance de la famille $(\sigma_i(X))_{i\geq0}$ des \uline{ens} $\sigma_i(X)$.
Pour la structure multiplicative, il faut connaître de plus les applications
\[
\sigma_i(X)\times\sigma_j(X)\longrightarrow
\mathbb Z^{\left(\coprod_{\underset{\sim}{k}\leq\unsure{i+j:2|i+j\cdot2|i+j-2}}\sigma_k(X)\right)}
\]
\inserted{$x,y\mapsto\deleted{\unsure{\xi}}uv$ défini par}
% archive sheet number: 39
\scanpage{41}
% L1
Nous allons expliciter plus bas la détermination de $\sigma^0(X)$ et $\sigma^1(X)$

% L3
Il y a lieu également de préciser les opérations $\vee$ et $\lambda^i$. Pour la première, notons que
\[
\deleted{\mathbb Z}\mathbb Q(-i)\,[M_i^+(X)]^\vee\subset M_i^+(X)
\]
d'où on conclut de proche en proche que $x\mapsto\overset{\vee}{x}\xi^i$ définit sur $\sigma_i$ une involution. Mais en fait il résultera de la théorie \unsure{projective} que \uline{cette dernière est l'identité}, i.e. \uline{l'opération} $\vee$ \uline{dans} \deleted{$\unsure{A}$} $M_i(X)$ \deleted{\unsure{est la dualité}} \uline{n'est autre que} $x\mapsto\xi^{-i}x$, \deleted{\unsure{d'où} $x\in\unsure{M_i}(X)\Longrightarrow\overset{\vee}{x}=\xi^{\unsure{-i|i}}x$}

% L14
Quant aux opérations $\lambda^i$, leur détermination devrait se faire en termes de groupes finis d'opérateurs (tels les $\unsure{\mathfrak S|G}_p$ \dots).
% archive sheet number: 40
\scanpage{42}
% L1
\uline{11}\quad \uline{Interprétation topologique du \deleted{\unsure{d}}\unsure{type|types}} \deleted{\unsure{filtré|filtr.}} \uline{dimensionnel} \unsure{(ou|(Cas} ``géométrique'').

% L3
Soit $X$ \deleted{(projective)} non singulière sur $k$ alg. clos. Considérons la filtration $X$ par la codimension, et la suite spectrale
\[
H^{\unsure{n|*}}(X,\mathbb Q_\ell)\Longleftarrow\inserted{E_1^{pq}=}\coprod_{x\in X[p]} H^{q-p}(x,\mathbb Q_\ell)(-p)\deleted{\unsure{)}}
\]
% slanted
\margin{Attention, terme initial $E_1$ et non $E_2$ !!}
\figure{\begin{tikzpicture}[x=0.7cm,y=0.7cm]
\draw (-0.65,0)--(3.9,0); \draw (0,-0.55)--(0,2.85);
\draw (0,0)--(4,2.9);
\node at (-0.32,2.95) {$q$}; \node at (4.8,-0.18) {$p$};
\node at (4.05,1.08) {$0$};
\end{tikzpicture}}
où on pose
\[
\left\{\begin{aligned}
X[p]&=\{x\in X\mid\dim\mathcal O_{X,x}=p\}\\
H^*(x,\mathbb Z_\ell(-p))&=\underset{\substack{U\text{ ouvert}\\\ne\varnothing\text{ de }\bar x}}{\underrightarrow{\lim}} H^*(U,\mathbb Q_\ell(\unsure{-p|p})))
\end{aligned}\right.
\]

% L11
On veut que \deleted{\unsure{la succession de}} cette suite spectrale provienne d'une suite spectrale de motifs [du moins à partir de $E_r^{pq}$ avec $r\geq2$ ; sinon il faudrait parler en ind-motifs, ou bien prendre une filtration \uline{finie} convenable de $X$ par des sous-schémas fermés]. Le morceau \inserted{en degré $h$} de la filtration $\geq p$ est visiblement de type dimensionnel $\leq h-2p$. On veut que \inserted{si $X$ est \uuline{projectif} (hypothèse essentielle \unsure{vu si|même si} $X$ de dim $1$ !)} \unsure{ce soient|ce sont|ce soit} \uline{exactement} les morceaux de type dimensionnel $h-p$.
% archive sheet number: 41
\scanpage{43}
% L1 — full passage crossed with long diagonal strokes
\deleted{si le \deleted{\unsure{fai}} motif de coefficient $\mathbb Z(0)$ est remplacé par \deleted{\unsure{le|un}} un motif constant tordu de \inserted{\unsure{purement|primitif}} type dimensionnel \deleted{\unsure{$d$}} $\delta$ [Attention au décalage par $\delta$ !]. Cet énoncé \deleted{\unsure{est}} \deleted{\unsure{de}} \uline{précise} les interprétations \unsure{galoisiennes|folkloriques|géométriques} et la conjecture de Tate du \deleted{\unsure{n°}}\dots, et est sans doute moins profond.}

% note crossing the boundary of the two sections
\margin{\deleted{$X$ de dim $n$, $U$ ouvert, $X-U$ de dim $<n$. Alors si $i\leq n$, $H^i(X)\deleted{\unsure{\simeq}}\longrightarrow H^i(U)$ est un isom \uuline{mod} $D_{i-1}(M^+_{\deleted{\unsure{i}}}(K))$, donc $H^i(X)$ et $H^i(U)$ ont $=$ composant ds $D_{i}(M^+(K))/D_{\unsure{n-1|i-1}}(M^+(K))$. De plus (via résolution) ce dernier est dans l'im de $M^+_{\unsure{i}}(K)$, d'où un inv. \unsure{birationel|birat.|birationnel} $\in\mathbb Z^{\sigma_i}$ \dots}}

% L9 — new heading
12) \uline{\deleted{\unsure{Bis|B.}}}\quad \uline{L'hom fondamental}
\[
L(K)\mathrel{\unsure{-|\longrightarrow}}M^+(K)
\]
\uline{et invariants birationnels fondamentaux}.

% L12
$K$ est un corps. On désigne par $L(K)$ le gpe $K$ engendré par les préschémas de type fini sur $K$, avec comme relations celles définies par découpage en morceaux. C'est un anneau par $\otimes$. On définit un hom d'anneaux
\[
\sigma:\quad L(K)\longrightarrow M^+(K)
\]
% archive sheet number: 42
\scanpage{44}
% L1
par
\[
\mathrm{cl}(X)\mapsto\sum(-1)^i\,\mathrm{cl}\,H^i_!(X,\mathbb Q(0))
\]
où bien entendu $H^i_!(X,\mathbb Z(0))$ est \deleted{\unsure{considéré}}\inserted{défini} comme motif sur $K$ par $R^if(\mathbb Q(0))$ ; $f:X\mathrel{\unsure{-|\to}}\operatorname{Spec}K$ la projection.

% L6
Prenant les composantes homogènes de $\sigma\mathrm{cl}(X)$, on a des $\sigma_i(\mathrm{cl}(X))$ qui sont des ``motifs virtuels'', et permettent de définir des ``nbs de Betti virtuels'', ceux demandés par Serre.

% L11 — diagonal crosses through the following paragraphs
\deleted{Si on filtre $L(K)$ par la dimension en des $L_i(K)$, on trouve}
\[
\deleted{L_i(K)\longrightarrow D_i(M^+_{\unsure{\xi|\varepsilon}}(K))}
\]
\deleted{On peut se demander si c'est surjectif, ou injectif, \deleted{ou bijectif mod} éventuellement mod groupes de torsion. Déjà pour $i=1$ ($K$ alg clos \inserted{de degré de trans. infini sur le corps premier}) rien n'est connu.}
% slanted
\margin{Pb pas bien posé \unsure{car avons|car on a} réponse \uline{négative} en \deleted{\unsure{telle}} \unsure{partie|partout}. cf plus bas la formulation correcte \dots\ En fait $L_i(K)\to M^+_{(i)}(K)+\xi M^+_{(i-1)}(K)+\xi^2 M^+_{(i-2)}(K)+\dotsb$}
\deleted{On peut corser la situation en considérant des $X/K$ munis d'un grpe fini donné $G$ d'autom. -- mais il faut}
% archive sheet number: 43
\scanpage{45}
% L1
Filtrons $L(K)$ par la dimension, on trouve
\[
L_{(i)}(K)\longrightarrow
\underbrace{M^+_{(i)}(K)+\underset{\substack{\cap\\ M^+_{(i+1)}(K)}}{\xi M^+_{(i-1)}(K)}
+\underset{\substack{\cap\\ M^+_{i+2}(K)}}{\xi^2 M^+_{i-2}(K)}+\dotsb}_{=D'_i(M^+(K))}
\]
\deleted{donc si $x\in$} \deleted{\unsure{De}}

% L5
d'où un hom. naturel
\[
\begin{aligned}
\mathrm{gr}_n(L(K))&\longrightarrow\mathrm{gr}'_{\deleted{\unsure{i}}n}(M^+(K))\\
&\hspace{2em}\big\Vert\quad\text{filtration par les }D'_i(M^+(K))\\
&\mathbb Z^{\deleted{\unsure{\coprod}}(\sigma_n\cup\xi\sigma_{n-1}\cup\xi^2\sigma_{n-2}\cup\dotsb)}
\end{aligned}
\]

% L9
Donc, si $X$ est un préschéma de type fini \deleted{de} sur $K$, de dim $\leq n$, alors prenant $\mathrm{cl}(X)\in L_n(K)$ et réduisant modulo $L_{n-1}(K)$, on trouve une image dans $\mathbb Z^{(\sigma_n\cup\xi\sigma_{n-1}\cup\dotsb)}$ donc un invariant de la forme
\[
\gamma_n(X)+\xi\gamma_{n-1}(X)+\xi^2\gamma_{n-2}(X)+\dotsb
\]
\[
\deleted{\beta_{\unsure{n}}(X)\unsure{+}\xi\beta_{n-1}(X)+\xi^2\beta_{n-2}\unsure{(X)}+\dotsb}
\]
\deleted{\uline{poids} \unsure{poids} $n\quad n+1\quad n+2$} [poids $n$, $n+1$, $n+2$, \dots]
% archive sheet number: 44
\scanpage{46}
% L1
où chaque
\[
\deleted{\unsure{\beta_{j}}}\gamma_{j}(X)\in\mathbb Z^{\sigma_j(X)}
\qquad 0\leq j\leq n
\]
Ils sont obtenus en prenant les $\beta_k(X)$ de Serre [composantes de poids purs de $\mathrm{cl}(X)$] pour $k\geq n$, et \deleted{les composant}\inserted{en notant} que un tel $\beta_k(X)$ se met sous la forme $\xi^{k-n}\beta'_{\deleted{\unsure{k}}2n-k}(X)$, où $\beta'_{\deleted{\unsure{k}}2n-\unsure{n|k}}(X)\in M^+_{2n-k}(K)$ ; alors $\gamma_{2n-\unsure{n|k}}(X)$ s'obtient en prenant la composante \uline{bipure} de $\beta'_{2n-k}(X)$ \dots\ .
% slanted
\margin{aux signes près !}

% L10
Par construction $=$, les invariants ainsi obtenus sont des invariants $n$-birationnels de $X$. Si $X$ est projective \deleted{non singulière}\inserted{lisse sur $K$}, alors $\beta'_j(X)$ pour $0\leq j\leq n$ n'est autre (par \uuline{Lefschetz} \dots) que la classe des $H^j(X,\mathbb Q(0))$. Donc notre construction consiste à prendre la composante $j$-bipure des $H^j(X,\mathbb Q(0))$. Mais cela n'est autre aussi
% archive sheet number: 45
\scanpage{47}
% L1
d'après le n\textsuperscript{o} précédent que \deleted{la composante \unsure{$j$-bipure|$j$-bipures} de} \inserted{l'image dans}
\[
\underset{\substack{U\,\text{ouvert}\\\deleted{\text{\unsure{de}}}\\\text{dense}}}{\underrightarrow{\lim}}H^j(U,\mathbb Q(0))
\]
de $H^j(X,\mathbb Z(0))$.

% L4
En fait, on constate (via dualité) que si $X$ est lisse (sans doute inutile) de dim $\leq n$, et si $i\leq n$ alors pour tt ouvert dense $U$ de $X$,
\[
H^i(X)\mathrel{\unsure{-|\to}}H^i(U)
\]
est \uline{\deleted{bijectif}\inserted{injectif} mod} $D_{i-2}(M^+(X))$ et \uline{bijectif mod} $\underline{D_{i-1}(M^+(X))}$, ce qui donne un invariant birationnel dans
\[
D_{i}(M^+(X))/D^{\unsure{+}}_{i-1}(M^+(X)),
\]
lequel est d'ailleurs dans l'image de $M^+_{(i)}(X)$, donc définit un objet de
\[
M^+_{(i)}(X)/[M^+_{(i-1)}(X)\unsure{;|)|,}\deleted{\unsure{+}}\xi M^+_{(i-2)}(X)],
\]
donc un élément de $\mathbb Z^{\sigma_i(X)}$. C'est l'invariant birationnel $\gamma_i(X)$.
% archive sheet number: 46
\scanpage{48}
% L1
Sous la forme \unsure{précédente|précédent}, on voit que l'on a mieux que des invariant dans $M^+_i(K)$. En fait, on a des
\[
\underline{\gamma}^{i}(X)\in\mathrm{Ob}\, M^+_{i\,\mathrm{bipure}}(K),
\]
objets semi-simples, définis de la façon suivante ($K$ parfait, ce qu'on peut supposer) $\gamma^i(X)\subset\underset{\substack{U\text{ ouvert}\\\text{dense}}}{\underrightarrow{\lim}}H^i(\unsure{U|M},\mathbb Z(0))$ est le plus grand sous-motif \uline{bipure} de type $i$. \deleted{[\unsure{Alors|la} le quotient par $\gamma^i(X)$} \inserted{C'est un foncteur en $X$} [où on l'anneau des fonctions rat. sur $X$, comme alg. sur $K$], de façon évidente
\diagram{\begin{tikzcd}[ampersand replacement=\&,row sep=small]
X'\arrow[d]\\X
\end{tikzcd}}
[NB si $X'\mathrel{\unsure{-|\to}}X$ n'est pas dominant, $X$ \unsure{irrid|irréd|irréd.}, alors $\gamma^i(X)\mathrel{-}\gamma^i(X')$ est nul]. Pour un hom \uline{dominant}, $\gamma^i(X)\longrightarrow\sigma^{i}(X')$ est \uline{injectif}. En effet, on peut supposer (prenant des modèles convenables) que $X'\mathrel{-}X$ est propre, \inserted{donc surjectif}, $X$ et $X'$ lisses, alors $H^i(X)\mathrel{-}H^i(X')$ est \uline{injectif} comme il est bien connu \dots

% L19
\uline{Question} \deleted{\unsure{la relation|Conditions}} Propriétés de ``fidélité'' du foncteur précédent ?? On aimerait que tout \deleted{hom} système d'hom $\underline{\deleted{\unsure{\beta}}\gamma}_i(\deleted{\unsure{K}}R)\mathrel{\unsure{-|\to}}\underline{\gamma}_i(R')$
% archive sheet number: 47
\scanpage{49}
% L1
provienne d'un $K$-hom $R\mathrel{\hookrightarrow}R'[t_1,\dotsb t_N]$ ($R,R'$ ext. de type fini de $K$), que l'isomorphie du système des $(\gamma_i(R))$ et du système des $(\gamma_i(R'))$ implique que $R$ et $R'$ sont ``équivalents'' en ce sens que chacun s'envoie dans une extension transcendante pure de l'autre \dots

% L9
\uline{Exemples} 1) Soient $C,C'$ deux courbes \inserted{lisses complètes \unsure{conn.|connex.|connexes}} \inserted{sur $k$ alg. clos,} supposons que les jacobiennes soient telles qu'il existe un hom $J\longrightarrow J'$ à noyau fini. Prouver qu'il existe \deleted{$C'$} une application rationnelle \unsure{dominant|dominante}
\[
C'[t_1,\dotsb t_N]\mathrel{\unsure{-|\to}}C\qquad??
\]
% slanted
\margin{\uline{$K$ alg. clos} \uline{au besoin} !}

% L15
2) Si les $\underline{\gamma}^{i}_{\deleted{\unsure{j|i}}}(R)$ sont nuls pour $i\unsure{>| \geq}0$, prouver que $R$ est unirationnelle ?? [En particulier, si la cohomologie \inserted{de $X$ var. proj. non sing.} est entièrement algébrique ??] Cas d'une surface ?? \inserted{[cl. alg. !]}
% archive sheet number: 48
\scanpage{50}
% L1
\uline{Remarque}. Si nous négligeons les \uline{variétés réglées} \inserted{de dim $n$}, [i.e. les extensions \inserted{$L$} de $K$ de \inserted{degré de tr. $n$} la forme $L_0[t]$, $t$ une indéterminée) le seul invariant \inserted{cohomologique} qui reste est $\underline{\gamma}_n$. Trouvons-nous un isom entre $\mathrm{gr}_n L(K)$ mod. le sous-grpe réglé [i.e. $\mathrm{Im}\,\xi\,\mathrm{gr}_{n-1}L(K)$] avec $\mathbb Z^{\sigma_n(K)}$ ? En particulier, si $\underline{\gamma}_n(X)\deleted{\text{\unsure{dire que}}}=0$, est-il vrai que $X$ est \uuline{réglée} ??

% L11
\uline{Exemple} Pour une surface : Si la cohomologie $H^2$ est algébrique, la surface est-elle réglée ??

\margin{\uline{Non en car. $p>0$}, prendre $E\times E$ avec $E$ d'invariant de Hasse nul.}
% archive sheet number: 49
